% TREŚĆ ROZDZIAŁU: zastąp wskazówki opisem testów i uzyskanych wyników.
% Każdą tabelę i każdy rysunek omów w tekście. Podawaj jednostki, warunki
% pomiaru oraz odwołanie do wymagania, które dany test weryfikuje.
\chapter{Weryfikacja i wyniki}
\label{chap:wyniki}

\section{Metodyka testów}
Opisz stanowisko, dane, warunki, procedurę i miary. Informacje powinny wystarczyć do powtórzenia najważniejszych testów.

\section{Wyniki}
Przedstaw wyniki w tabelach lub na wykresach, podając jednostki i warunki pomiaru. Nie powtarzaj w tekście wszystkich wartości widocznych w tabeli.

\section{Ocena spełnienia wymagań}
Dla każdego wymagania podaj wynik testu oraz status: spełnione, częściowo spełnione albo niespełnione. Odnieś się do kryteriów z \cref{chap:cel-wymagania}.

\section{Dyskusja}
Zinterpretuj rezultaty, wyjaśnij odstępstwa i wskaż ograniczenia wiarygodności oceny.
